\captionsenglish
\begin{resumo}
This dissertation presents the development and evaluation of an automated system for melanoma detection support in dermoscopic images, whose main methodological contributions are a five-axis framework for the automated assessment of structural and logical consistency in LLM-generated diagnostic reports, and an uncertainty calibration analysis revealing that the best-performing model concentrates its diagnostic hesitation precisely on the most clinically ambiguous cases. The architecture integrates expert-annotated segmentation masks (\textit{ground truth}) from the PH2 dataset with intelligent agents powered by multimodal large language models (LLMs), which apply three independent clinical protocols (ABCD, Menzies, and the Seven-Point Checklist (SPCL)) through a multi-agent approach implemented with the LangGraph framework. The PH2 dataset, comprising 200 expert-annotated dermoscopic images (40 melanomas and 160 non-melanoma lesions), served as the experimental base, with its segmentation masks adopted as standardized input to the system in order to isolate the diagnostic evaluation from the variability of automatic segmentation models. GPT-5.4-nano (OpenAI) and Gemini-3.1-flash-lite (Google) were evaluated and compared against a supervised ResNet-50 baseline for binary melanoma detection. Results show that the evaluated LLMs underperform the conventional baseline (ResNet-50, trained directly on PH2 with class-weighted loss: accuracy = 0.90; precision = 0.57; sensitivity = 1.00; F1 = 0.73, evaluated on a stratified test set) on most standard classification metrics, including sensitivity, an expected outcome for general-purpose multimodal models without domain-specific fine-tuning, but the five-axis consistency analysis uncovers a more informative pattern: Gemini-3.1-flash-lite outperforms GPT-5.4-nano across every axis, including ABCD arithmetic accuracy (100\% vs.\ 96\%), threshold adherence (46\% vs.\ 6\%), inter-protocol consensus (78\% vs.\ 46\%), and accuracy on definitive-diagnosis cases (85.2\% vs.\ 22.2\%). Critically, Gemini-3.1-flash-lite exhibits clinically coherent uncertainty calibration, concentrating its ``indeterminate'' and ``suspicious'' classifications on cases with higher real melanoma prevalence, whereas GPT-5.4-nano defaults to non-definitive diagnoses in 82\% of cases without discriminating risk levels. These findings position the five-axis consistency framework and the uncertainty calibration analysis as the core methodological contribution of this work, extending beyond raw classification accuracy toward a reproducible protocol for auditing the reliability of LLM-based diagnostic reasoning. The proposed architecture proves promising as an AI-assisted triage tool for melanoma detection, particularly in settings with limited access to dermatology specialists.

\textbf{Keywords: Melanoma detection, Uncertainty calibration, Diagnostic consistency, Multimodal language models, Multi-agent systems, PH2.}
\end{resumo}
\captionsbrazil
